\documentclass[11pt]{article}
% Code in a document: verbatim, lstlisting with a language, minted, and the
% inline forms \verb and \lstinline. Each becomes a code block (or an inline
% code run) in the imported document, with the language kept as the fence
% hint, so the renderer can highlight it.
\usepackage{listings}
\usepackage{minted}
\usepackage{xcolor}
\title{Code in Documents}
\author{The UltraCanvas Team}
\date{September 2026}
\begin{document}
\maketitle

\begin{abstract}
Verbatim material is the one place where a reader must stop reading LaTeX:
inside it, a backslash is a backslash. The importer switches the scanner off
for the whole body of these environments, keeps the lines exactly as they
stand, and records the language so the fence says
\texttt{```cpp} rather than \texttt{```}.
\end{abstract}

\section{Environments}
\begin{description}
  \item[verbatim] The plain one: no options, no language, every character
        kept.
  \item[lstlisting] Takes options; \texttt{language=} becomes the block's
        language hint.
  \item[minted] Takes the language as its argument:
        \texttt{\textbackslash begin\{minted\}\{python\}}.
\end{description}

\subsection{verbatim}
\begin{verbatim}
mkdir build && cd build
cmake .. -DULTRACANVAS_LATEX_ENGINE=native
make -j$(nproc)
\end{verbatim}

\subsection{lstlisting with a language}
\begin{lstlisting}[language=C++]
auto formula = CreateLaTeXView("eq1", 20, 20, 0, 0);
if (formula) {                       // null only if the module can't be loaded
    formula->SetTextSize(28.0f);
    formula->SetTextColor(Colors::Black);
    formula->SetLaTeX("\\frac{-b \\pm \\sqrt{b^2 - 4ac}}{2a}");
    window->AddChild(formula);
}
\end{lstlisting}

\subsection{minted with a language}
\begin{minted}{python}
def escape_velocity(mass, radius, G=6.674e-11):
    """Escape velocity in m/s, from the mass and radius of a body."""
    return (2 * G * mass / radius) ** 0.5

print(f"{escape_velocity(5.972e24, 6.371e6):.0f} m/s")
\end{minted}

\section{Code inside a sentence}
\texttt{\textbackslash verb} takes its delimiter from the character that
follows it, so
\verb|a_b^c| and \verb+x | y+ both come through untouched. The listings
package spells the same thing \lstinline{SetLaTeX}, and
\texttt{\textbackslash texttt} gives \texttt{monospaced text} that is not
code: the first two become code runs in the model, the third a plain run in
a monospaced family.

A code run survives every export: \lstinline{UCRichDocument} is a backtick
span in Markdown, a \texttt{<code>} element in HTML, and a character style
in ODT and DOCX.

\section{What is not kept}
Syntax highlighting is the renderer's business, not the importer's: colours
chosen by \texttt{\textbackslash lstset} or by a minted style are dropped,
and the language hint is passed on instead. Line numbers, frames and captions
attached to a listing go the same way --- the model stores code, and the
view decides how to show it.
\end{document}
